\section{Cyclic order-four central support}\label{sec:cyclic}
Write $K(C_4,4k)$ for the near-group ring
\[
 g^4=1,\qquad gX=Xg=X,\qquad
 X^2=1+g+g^2+g^3+4kX,
\]
where $k\ge1$ is an integer. The following reduction is uniform in $k$.
Thornton's theorem \cite[Theorem 1.3]{Thornton} permits a pseudo-unitary
Galois conjugate of any complex categorification of a near-group ring;
no braidedness or unitarity hypothesis is imposed on that reduction.
Set
\[
 d=2k+2\sqrt{k^2+1},\quad d^2=4+4kd,\quad
 D=4+d^2=8+4kd,\quad t=D/D'=1+kd=d^2/4.
\]
The dimension $d$ is irrational because $k^2<k^2+1<(k+1)^2$.

\subsection{A finite central-support lemma}
For a central simple $Z$ put
\[
 F(Z)=\sum_{r=0}^3m_{rZ}g^r+b_ZX,\qquad a_Z=\sum_r m_{rZ}.
\]
Call $Z$ \emph{mixed} if $a_Z>0$, and \emph{pure} otherwise.
The induction formula and adjunction give
\begin{equation}\label{eq:ng-gram}
 \sum_Zm_{rZ}m_{sZ}=4\delta_{rs}+1,\qquad
 \sum_Zm_{rZ}b_Z=4k,\qquad
 \sum_Zb_Z^2=8+16k^2.
\end{equation}

\begin{lemma}\label{lem:ng-support}
The mixed simples consist of six objects $B_{rs}$, $0\le r<s\le3$,
and two objects $E_{r,1},E_{r,2}$ for every $r$. Their restrictions are
\[
 F(B_{rs})=g^r+g^s+kX,\qquad
 F(E_{r,1})=g^r+u_rX,\quad F(E_{r,2})=g^r+(k-u_r)X,
\]
where $u_0=0$, $0\le u_r\le\lfloor k/2\rfloor$, and $u_1=u_3$.
For some $\eta\in\{i,-i\}$ their twists are
\begin{equation}\label{eq:ng-twists}
\begin{array}{c|c}
 E_{0,*},E_{2,*},B_{01},B_{02},B_{03}&1\\
 B_{12},B_{23}&-1\\
 B_{13}&\eta\\
 E_{1,*},E_{3,*}&-\eta.
\end{array}
\end{equation}
The set of $r$ for which $u_r=0$ is a subgroup of $C_4$. Its size
$h$ belongs to $\{1,2,4\}$ and equals the number of central invertibles.
Writing $L_r=u_r^2+(k-u_r)^2$, the total pure mass is
\begin{equation}\label{eq:ng-pure}
 \sum_{Z\,\mathrm{pure}}b_Z^2=8+10k^2-\sum_rL_r.
\end{equation}
If $h\ge2$, there is a central invertible $z$ with $F(z)=g^2$, $z^2=1$
and $\theta_z=1$. Splitting pure simples by their monodromy sign with
$z$ gives
\begin{align}
 \sum_{Z\in\mathcal P_+}b_Z^2&=4+6k^2-\sum_rL_r,
 &\sum_{Z\in\mathcal P_-}b_Z^2&=4+4k^2,\label{eq:ng-split}\\
 \sum_{Z\in\mathcal P_-}b_Z^2\theta_Z^j&=0
 &&\text{for every odd }j.\label{eq:ng-odd}
\end{align}
\end{lemma}
\begin{proof}
The formal codegrees are $D,D',4,4,4$. Thus
$I(1)=1\oplus A\oplus B_{01}\oplus B_{02}\oplus B_{03}$, with
$F(A)=1+kX$ and $F(B_{0r})=1+g^r+kX$; this follows from their dimensions
$1,t,D/4,D/4,D/4$ and \eqref{eq:ng-gram}. All have twist one.

The pointed associator class is trivial. Indeed, $\operatorname{add}(X)$
is a rank-one module category over $\C_{\mathrm{pt}}$; choosing the
isomorphisms $g^rX\simeq X$ gives a two-cochain whose coboundary is that
associator. Thus the pointed indicators equal those of $\Vec_{C_4}$:
$\nu_j(g^r)=1$ if $g^{rj}=1$ and zero otherwise. Positivity in
\eqref{eq:NS} at $j=\operatorname{ord}(g^r)$ forces
$m_{rZ}>0\Rightarrow\theta_Z^{\operatorname{ord}(g^r)}=1$.
All mixed twists therefore belong to $\{1,-1,i,-i\}$.

Apply the four-point Fourier transform to the traces for $I(g^r)$ at
$j=0,1,2,3$, with $j=0$ denoting dimension. Comparing coefficients of
$1,d$ is valid over $\Q(i)$ because $d\notin\Q(i)$. After summing over
$r$ the budgets are
\begin{equation}\label{eq:ng-budgets}
\begin{array}{c|rrrr}
 \theta&1&-1&i&-i\\\hline
 \sum a_Z^2&16&8&4&4\\
 \sum a_Zb_Z&8k&4k&2k&2k.
\end{array}
\end{equation}
Every $m_{rZ}$ is zero or one. For $r=0$ this follows from $I(1)$.
For $r\ne0$, $B_{0r}$ already contributes one to the diagonal Gram
entry five. An additional multiplicity two would consume the remaining
four; its overlap with either of the other two nonzero rows would then
be even, contrary to the off-diagonal entry one in \eqref{eq:ng-gram}.

A support of size three must avoid the unit row and therefore consist
of $g,g^2,g^3$. Its twist is $\pm1$ and it contributes nine to its
$a^2$-budget. But the negative budget is eight, while $I(1)$ already
contributes fourteen to the positive budget sixteen. A support of size
four is already impossible in the unit row. Each pair of rows
consequently overlaps in exactly one pair-support object, and each row
has two remaining single-support objects.

Neither $B_{12}$ nor $B_{23}$ fits in the remaining positive budget two,
so both have twist $-1$ and exhaust its budget. The two $g^2$-supported
objects must then have twist one. The object $B_{13}$ exhausts one
imaginary budget, leaving all four $g$- and $g^3$-supported singles with
the opposite imaginary twist. This proves \eqref{eq:ng-twists}.
The individual Fourier equations for $I(g)$ and $I(g^3)$ give mass
$D/4=2+kd$ in each fourth-root sector; those for $I(g^2)$ give $D/2$
in each sign sector. Comparing dimensions forces every pair coefficient
to be $k$ and every single pair to have sum $k$. Duality identifies the
unordered coefficient pairs over $g$ and $g^3$, so $u_1=u_3$.

Only a single-support object with coefficient zero has dimension one.
Forgetting embeds the group of central invertibles into $C_4$, because
the only invertible summand of $I(1)$ is the unit. This proves the
subgroup assertion and \eqref{eq:ng-pure}. If $h\ge2$, its order-two
object lies over $g^2$ and has twist one by \eqref{eq:ng-twists}.
Tensoring with it translates supports by two. All eight singles and
exactly $B_{02},B_{13}$ centralize it; the other four pairs have monodromy
$-1$. Their total mixed squared dimension is
\[
 16+16kd+(\sum_rL_r+2k^2)d^2.
\]
The centralizer has dimension $D^2/2$ by
\cite[Theorem 3.14]{DGNO}, and
\[
 D^2/2=16+16kd+(4+8k^2)d^2.
\]
Subtracting proves \eqref{eq:ng-split}. On the negative sector tensoring
with $z$ preserves $b_Z$ and changes the twist to $-\theta_Z$.
There are no fixed points, and pairing the terms proves
\eqref{eq:ng-odd}.
\end{proof}

\begin{lemma}[Balanced Galois obstruction]\label{lem:ng-galois}
A branch with $h=1$ and $u_1=u_2=u_3=k/2$ requires $k^2/8\in\ZZ$.
\end{lemma}
\begin{proof}
Here $A$ is the unique central simple of dimension $t$.
The signed Galois symmetry of the normalized modular matrix $S/D$
\cite[(2.13)]{BNRW}, at the embedding sending $d$ to $d'$, sends the
unit row to the $A$ row with positive sign. Hence
$S_{AZ}=t\,\sigma(\dim Z)$. Its eigenvalue ratios
$S_{AZ}/\dim Z$ are $t,t^{-1}$ at $1,A$, one at every other mixed
simple, and minus one at every pure simple; use $td'=-d$.
Since $A$ is self-dual, subtracting the unit fusion coefficient from
Verlinde's formula leaves only the two exceptional columns:
\[
 N_{AA}^Z-\delta_{Z,1}
 =\frac{(t^2-1)(\dim Z-\sigma(\dim Z))}{D^2}
 =\frac{k}{4}b_Z.
\]
The last equality uses $(t^2-1)(d-d')/D^2=k/4$.
Every single over a nonunit group element has $b_Z=k/2$.
Its fusion coefficient must therefore equal the integer $k^2/8$.
\end{proof}

\subsection{The parameterized indicator system}
Fix a moment height $H\ge3$. For each $u_1,u_2,u_3$ in Lemma~\ref{lem:ng-support} and
$\epsilon=\nu_2(X)\in\{-1,1\}$, take the known mixed data from
\eqref{eq:ng-twists}. Put
\[
 W_j=\sum_{Z\,\mathrm{mixed}} b_Z(a_Z+b_Zd)\theta_Z^j.
\]
Apply normalized relative trace to $K=\Q(\sqrt{k^2+1})$, interpreted
using its squarefree radicand. We may fix $\eta=i$, since the resulting
traces are unchanged by $i\mapsto-i$. Assign a nonnegative weight
$x_\tau$ to each pure joint trace type of Lemma~\ref{lem:trace-box}; it
represents the sum of $b_Z^2$ for that type. Its total is
\eqref{eq:ng-pure}. For $h\ge2$, instead use two pools with the masses
\eqref{eq:ng-split}. In the negative pool replace a joint tuple by
\[
 (\tau_j)_{j=1}^H\longmapsto
 \bigl(0\text{ for odd }j,\ \tau_j\text{ for even }j\bigr)_{j=1}^H,
\]
which is the average of the paired twists $\theta,-\theta$.
The first two equations are
\[
 \Tr^{\mathrm{norm}}_{L/K}(W_1)+d\sum x_\tau\tau_1=0,\qquad
 \Tr^{\mathrm{norm}}_{L/K}(W_2)+d\sum x_\tau\tau_2=\epsilon D.
\]
For orders $3\le j\le H$ with $F\nmid j$, impose the rationality and
bounds \eqref{eq:indicator-bound} on
$(\Tr^{\mathrm{norm}}_{L/K}(W_j)+d\sum x_\tau\tau_j)/D$.
The needed rotation dimensions are obtained without tensor powers:
if $X^j=\alpha_j\sum_rg^r+\beta_jX$, then
\[
 (\alpha_1,\beta_1)=(0,1),\qquad
 (\alpha_{j+1},\beta_{j+1})=(\beta_j,4\alpha_j+4k\beta_j),\qquad
 [X^j:1]=\alpha_j.
\]
The prime support is that of $2(k^2+1)$, since
$N_{K/\Q}(D)=64(k^2+1)$. All these systems are rational and become
integer systems after row scaling.

\begin{theorem}[Cyclic-four central-support criterion]\label{thm:cyclic}
Fix $k\ge1$ and a finite moment height $H\ge3$. Enumerate all the
coefficient and indicator branches above. If each branch is rejected by
Lemma~\ref{lem:ng-galois} or by an integer separating vector
\eqref{eq:dual-ray} for its central-support system, then $K(C_4,4k)$
admits no fusion categorification over $\CC$.
\end{theorem}
\begin{proof}
Any categorification has a pseudo-unitary Galois conjugate. Its central
support satisfies Lemma~\ref{lem:ng-support}, so belongs to an enumerated
branch. Lemma~\ref{lem:ng-galois} rules out the indicated balanced
branches. Every other branch supplies nonnegative weights in its
indicator system by \eqref{eq:NS}, the paired negative sector and
Lemma~\ref{lem:trace-box}. An integer separator contradicts their
existence exactly as in Theorem~\ref{thm:finite}.
\end{proof}

The software recognizes the entire cyclic-four tensor before applying
this criterion; recognizing only the number of invertible objects is
not sufficient. At height six it rejects the parameters $k=2,3,4,5$
encountered in the census beyond the positive case $k=1$. Their allocations are
reported together in Section~\ref{sec:automatic}; no separate numerical
root-of-unity argument is used for one selected ring.
